% page 185 of the facsimile (image p-184 of the scan)
% block 165.1 x 224.7 mm ; left margin 36.9 mm
\pgc{15.240mm}{0.000mm}{
\l{\cel{50}177}
\l{}
\l{\sou{fototipio}. Procedo di reprodukto fotomekanikala per imprim-inki,}
\l{\cel{3}en qua uzesas substanci koloida (gelatino, albumino, bitumo)}
\l{\cel{3}extensita sur fundi diversa (vitro, kupro, petro, zinko)}
\l{\cel{3}ed igita inkizebla per la lumo-radii. - DEF.}
\l{}
\l{\sou{fototropismo}. Efiko da la lumo-radii a la kresko di la planti.}
\l{}
\l{\sou{foxo}. (\sou{zool.}) Quadripedo karnavida, ek la genero \textquotedbl{}hundi\textquotedbl{}, di}
\l{\cel{3}qua la muzelo esas dina, la kaudo longa e tufa. - L.\cel{1}\sou{canis}}
\l{\cel{3}\sou{vulpes}. - DE.}
\l{}
\l{\sou{foxtroto}. Danso (zoologiala, ek Usa) quar-tempa, segun ritmo}
\l{\cel{3}inter marcho e polko. - E.}
\l{}
\l{\sou{foyo}. I.Singla de la kazi ye qui fakto eventas. - II. Singla}
\l{\cel{3}de la kazi ye qui quanto kontenesas kom elemento di toto. F.}
\l{}
\l{\sou{foyero}. (\sou{teatro)}. I. Loko ube la spektinti iris olim por varmi-}
\l{\cel{3}gar su dum la inter-akti, en la epoki dum qui la pleo-salono}
\l{\cel{3}ne esis kalorizata. - II. Salono, galerio, ube la spektinti}
\l{\cel{3}promenas dum la inter-akti. - DEFR.}
\l{}
\l{\sou{fraciono}. (\sou{matem.)}\cel{2}Un o plura parti di la uno dividita o subdi-}
\l{\cel{3}vidita a parti inter-egala. - DEFIRS.}
\l{}
\l{\sou{frago}. (\sou{bot.}) Frukto multa-bera, qua devenas de planto herbacea}
\l{\cel{3}de la familio \textquotedbl{}rozacei\textquotedbl{}. - L.\cel{1}\sou{fragaria}. - FISL.}
\l{}
\l{\sou{fragmento}. I. Peco de kozo ruptita o dislacerita. - II.(\sou{metaf.})}
\l{\cel{3}Io quo restas de libro di qua la maxim multo esas perdita o}
\l{\cel{3}permanis nefinita. Peco citita ek libro. - DEFIS.}
\l{}
\l{frajila. I. Ruptebla facile. - II. (\sou{metaf.}) Olqua facile cedas}
\l{\cel{3}a la tentesi. - DEFIS.}
\l{}
\l{frako. Viro-vesto nigra kun baski moruo-kaudatra, uzata por la}
\l{\cel{3}vesper-festi mondumala, la ceremonii. - DFRS.}
\l{}
\l{frakasar. (\sou{trans}.) Ruptar splitige. - FI.}
\l{}
\l{framo. Menuzuro quaza-kadra ; asemblajo de montanti e de traver-}
\l{\cel{3}si ; envelopilo buxatra. - E.}
\l{}
\l{framasono. Persono qua iniciesis en asociitaro (olim sekreta) qua}
\l{\cel{3}profesas principi di interfrateso, havas kom emblemi ula}
\l{\cel{3}instrumenti di masonisti, e di qua la membraro konsistas ye}
\l{\cel{3}grupi diversa (\textquotedbl{}logii\textquotedbl{}.) - EFIRS.}
\l{}
\l{frambo. (bot.)\cel{2}Frukto multa-bera, tre parfumoza, devenanta de}
\l{\cel{3}arboreto dornoza ek la familio \textquotedbl{}rozacei\textquotedbl{}. L.\cel{1}\sou{rubus idaeus.}FS.}
}
